% Construcción dodecafásica del funcional de retorno.
\section{Chaîne fondamentale dodécaphasique et fonctionnelle hexadique--octadique}

L’incidence active se construit avant toute réalisation gravitationnelle.
Soit
\[
 H_{\mathrm{act}}=\{1,2,4,5,7,8\},
 \qquad
 \mathcal P_2(H_{\mathrm{act}})
 =\{A\subset H_{\mathrm{act}}:|A|=2\}.
\]
Alors
\[
 |H_{\mathrm{act}}|=6,
 \qquad
 |\mathcal P_2(H_{\mathrm{act}})|=15,
\]
et les incidences antérieure et postérieure à l’extension ont pour cardinaux
\[
 \mathcal F_{90}=H_{\mathrm{act}}\times
 \mathcal P_2(H_{\mathrm{act}}),
 \qquad |\mathcal F_{90}|=90,
\]
\[
 \mathcal F_{120}=\{1,\ldots,8\}\times
 \mathcal P_2(H_{\mathrm{act}}),
 \qquad |\mathcal F_{120}|=120.
\]
Ces deux cardinaux sont des sorties de la règle de phase ; ce ne sont ni des angles ni
des paramètres reçus de la réalisation physique.

Notons \(C_{12}=\mathbb Z/12\mathbb Z\) le registre temporel du TPK et
\(\partial_{\rm ret}\) son unique couture de retour par tour. Dans le
module libre des événements orientés, on considère la chaîne fondamentale
\begin{equation}
 c_{12}^{+}
 =\sum_{j\in C_{12}}[j,\mathcal F_{90}]
  -[\partial_{\rm ret},\mathcal F_{120}].
 \label{intsuccpass:eq:c12-fundamental-chain}
\end{equation}
La somme contient exactement les douze secteurs de la roue ; le second
terme est unique parce qu’un tour possède une seule couture de clôture. Son signe
est le signe de frontière. L’involution bidirectionnelle envoie
\(c_{12}^{+}\) sur \(c_{12}^{-}=-c_{12}^{+}\) : elle change l’orientation de la
chaîne, mais n’altère ni ses multiplicités absolues ni ses types
d’incidence.

Pour \(m\in\{90,120\}\), soit
\[
 S_m(q)=\frac{q^m}{1-q^{3m}}.
\]
Le lecteur de retour \(\mathscr R\) est défini sur les générateurs de
\eqref{intsuccpass:eq:c12-fundamental-chain} par
\[
 \mathscr R([j,\mathcal F_{90}])=S_{90},
 \qquad
 \mathscr R([\partial_{\rm ret},\mathcal F_{120}])=S_{120}.
\]
La première égalité est constante sur \(C_{12}\) par covariance sous la
rotation de phase ; la seconde a un domaine distinct et n’agit que sur la
couture étendue. Par linéarité,
\begin{equation}
 \boxed{\mathscr R(c_{12}^{+})=12S_{90}-S_{120}.}
 \label{intsuccpass:eq:dodecaphase-forces-12minus1}
\end{equation}
Ainsi, les coefficients \(12:-1\) sont l’image de la chaîne fondamentale : douze
cellules de phase et une frontière orientée. Ils ne sont sélectionnés ni au moyen de la valeur
conventionnelle du paramètre de Barbero--Immirzi, ni au moyen d’une comparaison
décimale, ni en imposant au préalable le coefficient de pôle qui sera calculé
ci-après.

\begin{theorem}[Coefficient de pôle de la chaîne dodécaphasique]
La fonctionnelle \(F(q)=\mathscr R(c_{12}^{+})(q)
=12S_{90}(q)-S_{120}(q)\), déterminée par les douze secteurs
et leur couture orientée, a pour coefficient \(1/24\) dans
\((1-q)^{-1}\). Son résidu complexe en \(q=1\) est \(-1/24\).
Le coefficient de pôle est une sortie de la fonctionnelle déjà construite.
\end{theorem}

\begin{proof}
Pour une fonctionnelle ayant un pôle simple en \(q=1\), nous écrivons
\[
 p_1(F):=\lim_{q\uparrow1}(1-q)F(q).
\]
L’identité
\[
 1-q^{3m}=(1-q)\sum_{j=0}^{3m-1}q^j
\]
implique
\[
 p_1(S_m)
 =\lim_{q\uparrow1}\frac{q^m}{\sum_{j=0}^{3m-1}q^j}
 =\frac1{3m}.
\]
Par linéarité,
\[
 p_1\bigl(\mathscr R(c_{12}^{+})\bigr)
 =\frac{12}{270}-\frac1{360}
 =\frac{48-3}{1080}=\frac1{24}.
\]
Ces égalités s’effectuent dans \(\mathbb Q\). Le résidu complexe
est le coefficient de \((q-1)^{-1}\), de sorte que
\[
 \operatorname*{Res}_{q=1}\mathscr R(c_{12}^{+})(q)
 =-p_1\bigl(\mathscr R(c_{12}^{+})\bigr)=-\frac1{24}.
\]
\end{proof}

L’équation diophantienne
\[
 \frac a{270}-\frac b{360}=\frac1{24}
 \quad\Longleftrightarrow\quad
 4a-3b=45
\]
demeure un certificat arithmétique ultérieur. À elle seule, elle admet la
famille \((a,b)=(12+3t,1+4t)\) ; la chaîne
\eqref{intsuccpass:eq:c12-fundamental-chain} élimine cette ambiguïté avant le calcul,
car changer \((12,1)\) altère la multiplicité de l’orbite ou de sa couture
et, par conséquent, ne représente plus un tour fondamental du TPK.

Avec les canaux conjugués \(q_\pm\) déjà engendrés par la branche angulaire, on
obtient enfin
\[
 \mathscr R(c_{12}^{+})(q_-)-
 \mathscr R(c_{12}^{+})(q_+)
 =12\Delta S_{90}-\Delta S_{120}.
\]
C’est l’entrée interne de la fonctionnelle hexadique--octadique. Son évaluation
comme paramètre de Barbero--Immirzi et son insertion dans l’opérateur d’aire sont
des représentations ultérieures : elles ne participent à la sélection
ni de \(c_{12}^{+}\), ni de \(90/120\), ni de \(12:-1\), ni de \(1/24\).

\paragraph{Falsificateurs.}
La construction est réfutée si le recensement de phase ne donne pas douze secteurs, si
un tour contient plus d’une couture, si les incidences n’ont pas pour
cardinaux \(90\) et \(120\), si l’involution n’inverse pas la chaîne complète,
ou si le coefficient exact de \((1-q)^{-1}\) de son image diffère de \(1/24\).
